% Part: lambda-calculus
% Chapter: introduction
% Section: syntax

\documentclass[../../../include/open-logic-section]{subfiles}

\begin{document}

\olfileid{lam}{int}{syn}
\olsection{نحو حساب لامبدا}

کار را با دنباله‌ای از متغیرها چون~$x$، $y$، $z$،~\dots و چند نماد ثابت
چون~$a$، $b$، $c$،~\dots آغاز می‌کنیم. مجموعهٔ ترم‌ها به‌صورت استقرایی
چنین تعریف می‌شود:
\begin{enumerate}
\item هر متغیر یک ترم است.
\item هر ثابت یک ترم است.
\item اگر~$M$ و~$N$ ترم باشند، $(MN)$ نیز ترم است.
\item اگر~$M$ ترم و~$x$ متغیر باشد، آنگاه
  $(\lambd[x][M])$ نیز ترم است.
\end{enumerate}
ترم‌های به‌شکل~$(MN)$ را \emph{کاربرد} و ترم‌های به‌شکل
$(\lambd[x][M])$ را \emph{انتزاع} می‌نامند.

دستگاهی که هیچ ثابتی ندارد \emph{حساب لامبدای محض} نامیده می‌شود. ما
عمدتاً در حساب~$\lambd$ محض کار خواهیم کرد، پس همهٔ حروف کوچک نمایندهٔ
متغیرها خواهند بود. از حروف بزرگ ($M$، $N$ و مانند آن‌ها) برای نمایش
ترم‌های حساب~$\lambd$ استفاده می‌کنیم.

از چند قرارداد نمادگذاری پیروی خواهیم کرد:

\begin{conv}
\begin{enumerate}
\item وقتی پرانتزها حذف شوند، کاربرد از چپ به راست انجام می‌گیرد. برای
  نمونه، اگر~$M$، $N$، $P$ و~$Q$ ترم باشند، $MNPQ$ کوتاه‌نوشت
  $(((MN)P)Q)$ است.
\item باز هم وقتی پرانتزها حذف شوند، انتزاع لامبدایی گسترده‌ترین دامنهٔ
  ممکن را می‌گیرد. برای نمونه، $\lambd[x][MNP]$ به‌صورت
  $(\lambd[x][((MN)P)])$ خوانده می‌شود.
\item می‌توان با یک لامبدا چند متغیر را انتزاع کرد. برای نمونه،
  $\lambd[xyz][M]$ کوتاه‌نوشت
  $\lambd[x][\lambd[y][\lambd[z][M]]]$ است.
\end{enumerate}
\end{conv}

برای نمونه،
\[
\lambd[xy][xxyx \lambd[z][xz]]
\]
کوتاه‌نوشت عبارت زیر است:
\[
\lambd[x][\lambd[y][((((xx)y)x)(\lambd[z][(xz)]))]].
\]
باید این قراردادها را به خاطر بسپارید. در آغاز شما را کلافه خواهند کرد،
اما به آن‌ها عادت می‌کنید و پس از مدتی از سروکله‌زدن با انبوهی از
پرانتزها کمتر کلافه‌کننده خواهند بود.

دو ترمی که تنها در نام متغیرهای مقید تفاوت دارند، $\alpha$-هم‌ارز
نامیده می‌شوند؛ برای نمونه، $\lambd[x][x]$ و~$\lambd[y][y]$. مناسب است
آن‌ها را ترمی «یکسان» بدانیم؛ به بیان دیگر، وقتی می‌گوییم~$M$ و~$N$
یکسان‌اند، منظورمان «تا تغییر نام متغیرهای مقید» نیز هست. متغیرهایی که
در دامنهٔ یک~$\lambd$ قرار دارند «مقید» و بقیه «آزاد» نامیده می‌شوند.
در نمونهٔ پیشین هیچ متغیر آزادی وجود ندارد؛ اما در
\[
(\lambd[z][yz])x
\]
$y$ و~$x$ آزاد و~$z$ مقیدند.

\end{document}
